% ITSC 2026 Submission
% Hierarchical MARL for Distributed-Power x Electrical/Pneumatic Braking Division
% of Heavy-Haul Trains on Long Downhill Sections
% IEEE Conference Template (ITSC 6-page format)
%
% ARCHITECTURAL REWRITE (ADR-009 / 2026-09-05): replaces the previous
% "MARL vs MPC / threshold-coordination" draft. Core = per-locomotive ED/AB
% division MARL. All numbers from DIAGNOSTICS.md 2026-09-05 decompose finals
% (interval=10, 15,000 steps). Aligned with ITSC_PAPER.md.

\documentclass[conference]{IEEEtran}
\IEEEoverridecommandlockouts

\usepackage{amsmath,amssymb,amsfonts}
\usepackage{graphicx}
\usepackage{xcolor,booktabs,multirow}
\usepackage{makecell}
\newcommand{\permille}{\textperthousand}

\begin{document}

\title{Hierarchical MARL for Distributed-Power $\times$ Electrical/Pneumatic
Braking Division of Heavy-Haul Trains on Long Downhill Sections}

\author{
\IEEEauthorblockN{Anonymous Author}
\IEEEauthorblockA{\textit{Anonymous Institution} \\
City, Country \\
email@example.com}
}

\maketitle

%-------------------------- ABSTRACT --------------------------
\begin{abstract}
Heavy-haul trains on sustained long downhill sections must hold speed within a
narrow band using \emph{cyclic braking} (electric $+$ pneumatic), while keeping
in-train \emph{coupler forces} below a safety limit --- a long-horizon,
switched, safety-critical control problem. We show that the \emph{decomposition}
of this problem determines whether multi-agent reinforcement learning (MARL)
can learn anything at all: prior ``train-level'' MARL treated each whole train
as an agent and only nudged a shared speed-band threshold (a $\pm 0.5$~km/h
offset $=$ 1.25\% of bandwidth), leaving a lever too small to learn from (a
``nothing-to-learn'' ceiling). We re-decompose at the \emph{locomotive} level:
each locomotive is an independent agent that directly modulates its own
\emph{electric-braking depth} and \emph{pneumatic-braking trigger} --- a
continuous, physically meaningful lever. A rule engine (ED-first $+$
inter-locomotive stagger $+$ AB fallback) provides a safety-embedded inner
layer; MAPPO adds per-locomotive $[\mathrm{ed}_{mod},\mathrm{ab}_{mod}]
\in[-1,1]^2$ modulation. A semi-MDP \emph{decision interval} splits the
15{,}000-step credit chain into 1{,}500 decisions, making long-horizon MAPPO
converge. On the real Shuozhou--Huanghua profile (23.8~km, $\approx
-10.9$\permille) with a 10{,}000-t, 1+1 distributed-power consist, we obtain a
three-layer positive evidence chain with zero safety violations:
single-locomotive rule ($+40{,}956$) $\rightarrow$ distributed-power rule
division ($+71{,}252$; $+74\%$ reward, AB usage $-89\%$) $\rightarrow$ RL
division ($+72{,}743$; coupler peak $-21\%$ to \textbf{570~kN}, final speed
$+11.2$~km/h, ATP emergency events $-97\%$). The learned policy is physically
interpretable --- halve ED depth and inject staggered \emph{minimum-notch} AB
to cut longitudinal impulse --- reaching a region no tuning of the old
threshold coordination could attain.
\end{abstract}

\begin{IEEEkeywords}
Heavy-haul train, distributed power, electro-pneumatic braking division,
multi-agent reinforcement learning, cyclic braking, long downhill, coupler
force.
\end{IEEEkeywords}

%-------------------------- I. INTRODUCTION --------------------------
\section{Introduction}
\label{sec:intro}

A 10{,}000-ton heavy-haul train on a long downhill section (Shuozhou--Huanghua
line, 23.8~km, sustained up to $-10.9$\permille) must regulate speed through
\emph{cyclic braking}: brake above an upper speed limit, release below a lower
limit. The minute-scale brake--release cycles, combined with multi-body coupler
dynamics and finite air-brake wave propagation ($\approx 200$~m/s across a
$\sim 4$~km train), generate in-train \emph{coupler forces} that must stay
below 1000~kN (normal) / 2250~kN (emergency) \cite{fu2019}. Electric braking
(ED) is smooth, continuous, and available only on locomotives; pneumatic (air)
braking (AB) is discrete (RELEASE/MIN/SERVICE/EMERGENCY) and is the dominant
source of longitudinal impulse while simultaneously being slow to recover (pipe
recharge). How to \emph{split} braking effort between ED and AB across multiple
locomotives --- the \emph{distributed-power (DP) $\times$ electro-pneumatic
division} problem --- is exactly where controller design matters most on long
downhills \cite{tang2020,zhao2022}.

Our first attempt at MARL for this regime followed the literature's
``train-level'' decomposition \cite{liu2025}: one agent per \emph{train},
outputting a small offset to a shared speed-band threshold. The result was a
hard ceiling. The offset is $\pm 0.5$~km/h against a 40--80~km/h band --- a
\textbf{1.25\% bandwidth lever} --- and across a 15{,}000-step run its reward
gradient is statistically invisible. Adding global observations and
consistency-shaped rewards recovered only $\approx 0.8\%$ ($\sim 343$) reward.
We diagnosed this as a \emph{structural} problem, not a tuning problem:
\textbf{MARL has nothing to learn when its only control lever is too small to
change the outcome.}

The resolution is to re-decompose the action space at the \emph{locomotive}
level and give MARL a \emph{real} physical lever:
\begin{itemize}
\item \textbf{Electric-braking depth} --- each locomotive's $ED\in[0,1]$ is
continuous and directly scales the applied ED force (\emph{big} lever, $\pm 0.5$
of absolute ED depth).
\item \textbf{Air-brake triggering} --- each locomotive may suppress or
pre-trigger its own AB source along the consist (inter-locomotive
\emph{stagger} of pneumatic application), the primary knob for coupler impulse.
\end{itemize}

A rule engine provides the safety-embedded inner layer: ED-first,
per-locomotive stagger, AB only as a fallback when ED availability is
power-starved. MAPPO then learns how to \emph{modulate} this baseline per
locomotive. Because both ED depth and the AB on/off timing are physically
meaningful, the RL policy can move coupler force, braking energy, and holding
accuracy together.

A second, orthogonal obstacle is \textbf{credit assignment over a very long
horizon}. The full 23.8~km section is a 15{,}000-step time domain;
single-step-decision MAPPO over that horizon produced a flat, negative training
reward (the RL policy ended pervasively dragging the train with air braking).
We cure this with a \textbf{semi-MDP decision interval}: the RL policy decides
once every $N{=}10$ steps, holds its action, and stores a window-summed
transition --- compressing the 15{,}000-step credit chain to 1{,}500 decisions
while the rule baseline still executes per step. This is the general technique
on which the headline results depend.

\textbf{Contributions.}
(C1) We identify and formalize the ``\emph{nothing-to-learn}'' ceiling of
train-level threshold coordination (a 1.25\% bandwidth lever) and show that
re-decomposition to the locomotive ED/AB level provides a physically meaningful
MARL lever.
(C2) We give a \emph{three-layer, like-for-like positive evidence chain} on the
real Shuozhou profile --- single-locomotive $\rightarrow$ distributed-power rule
division $\rightarrow$ distributed-power RL division --- all with zero safety
violations, isolating the \emph{architecture} gain from the \emph{learning}
gain.
(C3) We establish \emph{semi-MDP decision intervals} as a practical cure for
long-horizon (15{,}000-step) MARL credit assignment in switched train dynamics.
(C4) The learned policy is \emph{physically interpretable} --- halve ED depth
$+$ staggered MIN-notch AB --- and cuts coupler peak by 21\%, raises final speed
by 11.2~km/h, and reduces ATP-triggered emergency events by 97\%.

%-------------------------- II. RELATED WORK --------------------------
\section{Related Work}
\label{sec:related}

\textbf{Distributed power $\times$ braking division.} Tang et al.\
\cite{tang2020} model multi-locomotive trains with RL and explicitly include
coupler force in the reward, establishing the DP setting we build on. Zhao et
al.\ \cite{zhao2022} give a refined air-brake model and report, for a 1+1
distributed-power consist, that \emph{synchronous multi-source braking reduces
in-train impulse} (their two-source $\approx 616$~kN vs.\ single-source
$\approx 1549$~kN in our re-derivation), directly supporting the ``ED/AB split
$+$ stagger'' lever. Our contribution is to make this division \emph{learned}
per locomotive via MARL rather than hand-tuned.

\textbf{RL for cyclic braking on long downhills.} Zhang et al.\ \cite{zhang2024},
Liu et al.\ \cite{liu2021}, and Wang et al.\ \cite{wang2025} all apply
DQN-type RL to long-downhill automatic speed control --- and \emph{all
explicitly ignore coupler forces}. We retain their cyclic-braking formulation
but add coupler force as a first-class objective/constraint (1000~kN), which is
precisely where the ED/AB division matters (AB is the impulse source). Zhang et
al.\ \cite{zhang2t2021} and Yang et al.\ \cite{yangmin} document that coupler
force stays below 1000~kN on mild grades when braking is released
simultaneously and pipe recharge is sufficient --- the physics our dynamic
multi-body model reproduces. Fan et al.\ \cite{fan2024} quantify how consist
marshalling shapes longitudinal impulse in cyclic braking.

\textbf{Why train-level MARL has ``nothing to learn.''} Prior group-platoon
MARL (e.g., the multi-train formulation in the earlier draft of this work)
gave each \emph{train} a cooperative threshold offset. The offset is $\pm 0.5$
km/h on a 40-km/h band $\rightarrow$ a 1.25\% lever $\rightarrow$ near-zero
reward gradient (our diagnostics). We report this as a \emph{structural
caution} for MARL action design: the lever must be able to change the outcome,
or the gradient vanishes regardless of observation/architecture
\cite{basile2022,liu2025}.

\textbf{Long-horizon credit assignment.} PPO/MAPPO over tens of thousands of
steps is known to degrade when the reward that matters arrives after many
decisions. Our semi-MDP decision interval (action hold $+$ window return) is a
minimal, task-agnostic mitigation that preserves the underlying dynamics while
shortening the credit chain; we include it as a reusable technique.

%-------------------------- III. PROBLEM & MODEL --------------------------
\section{Problem Statement and System Model}
\label{sec:model}

Longitudinal dynamics per vehicle (NeTrainSim-derived multi-body integration):
traction $-$ electric braking $-$ pneumatic braking $-$ Davis resistance
$\pm$ grade $-$ coupler force. Air braking propagates along the consist at
finite wave speed ($\approx 200$~m/s $\rightarrow$ seconds of delay across a
$\sim 4$~km train); electric braking is continuous and limited by power and
adhesion. Coupler forces come from a mass--spring--damper with a
\emph{gap dead-zone} (0.01~m), reproducing the ``coupler gap $\rightarrow$
coupler force'' law \cite{yangmin} and Newton-pair cancellation.

Control objective: keep each locomotive's speed in $[v_{\min}, v_{\max}] =
[40, 80]$~km/h, keep coupler force $\le 1000$~kN, over the full 23.8~km
(15{,}000-step) horizon, using a 1+1 distributed-power consist (10{,}000~t).
The difficulty is that the minute-scale cyclic rhythm and the band-limited
coupler constraint act over far longer than a rolling window, and the AB
actuator is discrete $+$ history-dependent (recharge).

\textbf{Why ED/AB splitting resists a small-lever policy.} On this profile the
rule-based cyclic braking already holds the band; the \emph{remaining} degrees
of freedom --- exactly how deep each locomotive brakes with ED, and exactly
when each triggers AB --- are the ones that decide coupler impulse. A policy
that cannot touch these (e.g., only shifts the band) is orthogonal to the
objective it must minimize.

%-------------------------- IV. PROPOSED METHOD --------------------------
\section{Proposed Method: Two-Level Decompose MARL}
\label{sec:method}

\textbf{Level 1 --- rule engine (safety-embedded inner layer).}
A deterministic baseline \texttt{get\_action\_decomposed} runs per step:
\emph{ED-first} (both locomotives saturate ED up to the physically available
share), \emph{per-locomotive stagger} (offset each locomotive's brake phase to
spread impulse), and \emph{AB fallback} only when ED availability is
insufficient. Critical to correct triggering is a \emph{physical ED-availability
model}:
\begin{equation}
F_{ED}=\min\!\big(P_{ED}/v,\; F_{\text{adhesion}}\big),\qquad
F_{\text{req}} \propto m\,g\,\max|\text{grade}|.
\end{equation}
On a steep section at speed the available ED share $<1$, so the baseline
genuinely falls back to AB; earlier simplifications ($\mathrm{clip}(v/5)$) made
the fallback unreachable and hid the mechanism.

\textbf{Level 2 --- MARL modulation (learned).} MAPPO, \emph{one agent per
locomotive} (1+1 consist $\rightarrow$ 2 agents), observes per-locomotive state
and outputs $[ed_{mod}, ab_{mod}]\in[-1,1]^2$:
\begin{itemize}
\item $ed_{mod}$ directly scales this locomotive's ED depth ($\pm 0.5$
absolute) --- the continuous physical lever;
\item $ab_{mod} > +0.5$ suppresses this locomotive's AB source (only while
$v < v_{\max}-0.5$, so the safety fallback is never cheaply switched off);
$ab_{mod} < -0.5$ pre-triggers a MIN-notch AB --- enabling \emph{staggered,
anticipatory} pneumatic braking.
\end{itemize}

\textbf{Observation (11-dim per locomotive).} Speed error w.r.t.\ target;
margin to $v_{\max}$; margin to $v_{\min}$; the rule baseline's ED (what the
engine ``wants''); whether the rule baseline would brake with AB; local coupler
force (normalized); AB-active fraction on the consist; AB release-remaining
fraction (recharge state); brake-segment index; locomotive position (for grade
inference); inter-train gap error (0 for single-train). Feeding the
\emph{air-brake state} lets the policy \emph{sense} hysteresis rather than
predict it.

\textbf{Reward.} Band holding (in-band $+5$, out-of-band $-20$ per train) $+$
coupler force (mean \& peak penalties, plus a super-limit term) $+$ comfort $+$
gap (multi-train extension) $+$ an ED micro-cost ($-0.02$ per depth) to avoid
gratuitous full-ED cruising. Notably, \textbf{AB is not directly penalized}:
extensive smoke-testing showed any positive AB cost exactly cancels the
coupler-force benefit of staggered MIN-notch AB (coupler force is the
objective; AB is only a means).

\textbf{Semi-MDP decision interval.} With a per-step decision and a 15{,}000-step
horizon, MAPPO's stage-3 training reward stayed flat at $\approx -291$~k ---
the credit chain was too long. With decision interval $N{=}10$, the RL policy
decides every 10 steps, holds its action, and stores one transition with a
window-summed return; the rule baseline still executes every step. This turned
a failed run into the converged results below.

%-------------------------- V. EXPERIMENTS --------------------------
\section{Experiments}
\label{sec:experiments}

\textbf{Protocol.} Shuozhou--Huanghua main profile (23.8~km), 1+1 distributed
power (10{,}000~t), ETCS distance-to-go curves on, deterministic evaluation
(fixed seed) at 80~km/h over the full 15{,}000 steps, $n_{\text{seed}}=3$
evaluation rollouts. MAPPO: $\mathrm{lr}=5\times10^{-5}$, entropy
$=0.002$, curriculum $1000{:}0,\,3000{:}20,\,15000{:}40$, decision interval
$N{=}10$, 60 episodes. The same environment $+$ same rule baseline are used for
M0 and the single-locomotive runs --- the three rows of
Table~\ref{tab:main} differ only in the \emph{architecture/structure} under
test (like-for-like).

\begin{table*}[!t]
\centering
\caption{Main result --- three-layer evidence chain (15{,}000 steps, all zero
safety violations).}
\label{tab:main}
\begin{tabular}{@{}l r r c r r r r r@{}}
\toprule
Method & reward & \makecell[c]{coupler peak\\(kN)} & viol & \makecell[c]{band\\rate} & \makecell[c]{AB\\steps} & \makecell[c]{mean\\ED} & \makecell[c]{final speed\\(km/h)} & \makecell[c]{ATP\\events}\\
\midrule
\textbf{decompose-M1 (RL division, $N{=}10$, ep60)} & \textbf{$+72{,}743$} & \textbf{570} & 0 & 99.72\% & 4{,}357 & 0.35 & \textbf{52.9} & \textbf{1.44\,M}\\
decompose-M0 (distributed-power rule division) & $+71{,}252$ & 720 & 0 & 99.68\% & 461 & 0.67 & 41.7 & 43.2\,M\\
single-locomotive (old 1+0 rule) & $+40{,}956$ & 456 & 0 & 90.20\% & 4{,}030 & --- & 44.9 & 467\,M\\
\bottomrule
\end{tabular}
\end{table*}

\textbf{Isolated increments.}
(a) \emph{Architecture gain} (M0 vs.\ single-loco): distributed-power
ED-first $+$ stagger $+$ ED-availability-faithful AB fallback gives
\textbf{$+74\%$ reward} ($+30{,}278$), $+9.5$~pp band rate, and \textbf{$-89\%$
AB usage} (4{,}030$\rightarrow$461 steps) --- distributed ED power suffices, so
pneumatic effort almost vanishes, with zero super-limit.
(b) \emph{Learning gain} (M1 vs.\ M0): RL finds a \emph{physically
interpretable re-allocation} --- halve ED depth (mean ED $0.67\rightarrow
0.35$) and inject \emph{staggered MIN-notch AB} (461$\rightarrow$4{,}357 steps)
--- trading slightly more, but individually gentle and phased, pneumatic
applications to shrink longitudinal impulse: \textbf{coupler peak $-21\%$
(720$\rightarrow$570~kN)}, \textbf{final speed $+11.2$~km/h}, \textbf{ATP
emergency events $-97\%$} (43.2\,M$\rightarrow$1.44\,M). The smoother ED
profile plus phased AB is exactly the region a small-lever threshold policy can
never reach.

\textbf{Why this is the decisive comparison.} The three rows share the
environment, the safety layer, and the rule baseline; only whether (i) power is
distributed and (ii) ED/AB are RL-modulated per locomotive changes. Every layer
adds a positive, \emph{measured} increment with zero violations ---
decomposing the source of the gain (architecture vs.\ learning) rather than
reporting a single black-box number. The same simulator, given a real lever,
transitions from a flat $-291$~k training reward to $+72.7$~k rollout reward.

\textbf{Ablation --- lever isolation (which lever carries the gain).} We reran
full training (seed0, same protocol) with one modulation masked during
training (\texttt{mask\_lever}), isolating the causal contribution of each
lever (best-checkpoint evaluation; all zero violations, band $\approx$ 100\%):
Table~\ref{tab:ablation}.

\begin{table}[!t]
\centering
\caption{Lever-isolation ablation (seed0 best checkpoint).}
\label{tab:ablation}
\begin{tabular}{@{}l r r r r r@{}}
\toprule
Method & reward & \makecell[c]{coupler\\(kN)} & \makecell[c]{AB\\steps} & \makecell[c]{mean\\ED} & \makecell[c]{final\\(km/h)}\\
\midrule
\textbf{M1 both levers (seed0 best)} & $+73{,}733$ & 578 & 4{,}807 & 0.34 & 53.5\\
only-ED lever (mask ab) & $+71{,}932$ & \textbf{528} & \textbf{0} & 0.70 & 45.1\\
only-AB lever (mask ed) & \textbf{$+74{,}634$} & 702 & \textbf{5{,}960} & 0.05 & \textbf{71.0}\\
M0 rule (no levers) & $+71{,}252$ & 720 & 461 & 0.67 & 41.7\\
\bottomrule
\end{tabular}
\end{table}

The two levers serve distinct objectives. The \textbf{ED lever is the
coupler-safety lever}: ED-only reaches the \emph{lowest} coupler peak
(\textbf{528~kN}, 27\% below M0) with zero AB and 100\% band, so the coupler
reduction of M1 vs.\ M0 comes chiefly from ED-depth shaping, not from AB. The
\textbf{AB lever is the throughput lever}: AB-only rides the speed-band ceiling
(final \textbf{71.0~km/h}, heavy phased AB $\approx$5{,}960 steps, near-zero
ED) and scores the \emph{highest} scalar reward ($+74{,}634$), at the cost of
higher coupler force (702~kN). M1 (both levers) is the compromise --- 578~kN,
$+73{,}733$, final 53.5~km/h --- jointly minimizing coupler force and holding
throughput. \textbf{Honest caveat:} only-AB outperforms M1 on the raw scalar
reward, i.e., the current reward shaping rewards high-speed band-riding
air-braking; we therefore do \emph{not} claim M1 is the scalar-reward optimum,
but the best \emph{coupler--throughput balance} among safe (zero-violation)
policies. Reshaping to penalize heavy AB (impulse/energy, not a linear cost) is
a stated follow-up.

\textbf{Generalization / recentring.} Because the policy acts on
per-locomotive ED/AB (not on a train-wide band), it transfers structurally to
longer sections and larger numbers of distributed-power locomotives without
re-tuning the band; ablations over $N\in\{5,10,20\}$ and AB-cost
$\in\{0,0.05\}$ (planned) will establish the margins of the semi-MDP interval
and the reward shape.

%-------------------------- VI. CONCLUSION --------------------------
\section{Conclusion and Limitations}
\label{sec:conclusion}

We argued and demonstrated that for long-downhill cyclic-braking control,
\textbf{the decomposition of the MARL action space decides whether there is
anything to learn}. Train-level threshold coordination gives a 1.25\%
bandwidth lever and a near-zero gradient (``nothing to learn'');
locomotive-level ED/AB division gives continuous, physically meaningful levers.
Combining a safety-embedded rule engine with per-locomotive MAPPO modulation
and a semi-MDP decision interval yields a three-layer, zero-violation positive
evidence chain on the real Shuozhou profile (RL: $+2.1\%$ reward over rule
division; coupler $-21\%$; final speed $+11.2$~km/h; ATP events $-97\%$).

\textbf{Limitations.}
(1) The learned policy is trained in a high-fidelity simulator and inherits its
fidelity (a sim-to-real gap shared in kind with model-driven MPC).
(2) We provide no formal recursive-stability guarantee; deployment safety would
require formal verification / ATP integration.
(3) AB is not directly penalized in the reward --- the division is only as safe
as the coupler-force weighting, and the AB-trigger margin
($v_{\max}-0.5$) is a hand-set constant.
(4) Long-horizon training remains expensive and seed variance is nontrivial;
the semi-MDP interval $N$ is a hyperparameter worth ablating.
(5) Inter-train gap adaptation (platoon level) is a stated extension; this
paper fixes a single train with 1+1 distributed power.
(6) The spring-buffer (draft-gear) model introduces a few calibrated constants:
the MT-2 static/dynamic friction ratio $\mu_s/\mu_d=1.5$ and transition velocity
$v_{\mathrm{er}}=0.2$~m/s have a literature-anchored structure (Yang Liangliang
2017, Eq.~3) but point values from impact tests; the buffer is unified to a
single MT-2 friction type (including the locomotive--freight connection),
avoiding an elastic-paste sub-model whose five shape constants had no published
fits;
our Shuozhou-profile sweep shows the coupler-force peaks shift $\le$1\%
within reasonable uncertainty of these values
($v_{\mathrm{er}}$ 0.05--0.5~$\to$~+1\%; $\mu_s/\mu_d$ 1.2--1.5~$\to$~+1.2\%),
so they do not dominate the headline coupler-force result (though $\mu_s/\mu_d$
shows a threshold nonlinearity near 1.8). Also, the release-time constant
$\tau=6$~s models only the initial exhaust transient (exponential decay); full
air-brake recharge takes 175--512~s (Zhao~2022), which we abstract as a
history-dependent recharge effect rather than simulate explicitly.

\begin{thebibliography}{1}

\bibitem{tang2020} G.~Tang, H.~Wang, and Y.~Liu, ``Reinforcement learning
approach for optimal control of multiple electric locomotives in a heavy-haul
freight train: A double-switch Q-network architecture,'' \emph{Knowledge-Based
Systems}, 2020.

\bibitem{zhao2022} Z.~Zhao, ``Refined modeling and operation optimization of
air braking for heavy-haul trains'' (in Chinese), M.S. thesis, Southwest
Jiaotong University, Chengdu, China, 2022.

\bibitem{zhang2024} J.~Zhang, X.~Zhou, Z.~He, \emph{et al.}, ``Deep Q-network
algorithm-based cyclic air braking strategy for heavy-haul trains,''
\emph{Algorithms}, vol.~17, no.~5, p.~190, 2024.

\bibitem{liu2021} S.~Liu \emph{et al.}, ``DDQN-based cyclic braking control of
heavy-haul trains on long downhill sections,'' \emph{Transportation Research
Part C}, 2021.

\bibitem{wang2025} Z.~Wang \emph{et al.}, ``DQN-based speed control of
heavy-haul trains on long downhill sections with cyclic-braking constraints,''
in \emph{Proc.\ IEEE Int.\ Conf.\ Intelligent Rail Transportation (ICIRT)},
2025.

\bibitem{zhang2t2021} E.~Zhang \emph{et al.}, ``Analysis and
application of longitudinal coupler forces of 20{,}000-t heavy-haul combined
trains on the Shuozhou--Huanghua railway'' (in Chinese), \emph{Shenhua
Science and Technology}, vol.~19, no.~1, pp.~74--78, 2021.

\bibitem{yangmin} M.~Yang \emph{et al.}, ``Influence of coupler gap on
longitudinal impact dynamics of heavy-haul trains'' (in Chinese), CNKI.

\bibitem{fan2024} Y.~Fan \emph{et al.}, ``Influence of consist
marshalling on longitudinal impulse of a 30{,}000-t heavy-haul train under
cyclic braking'' (in Chinese), \emph{Electric Locomotives \& Mass Transit
Vehicles}, no.~1, 2024.

\bibitem{fu2019} Y.~Fu, J.~Yuan, Z.~Li, and
H.~Yang, ``Driving process optimization of heavy-haul trains based on
coupler-and-buffer constraints'' (in Chinese), \emph{Acta Automatica Sinica},
vol.~45, no.~12, pp.~2355--2365, 2019.

\bibitem{wu2023} Q.~Wu \emph{et al.}, ``Railway virtual coupling: a survey of
emerging control techniques,'' \emph{IEEE Trans.\ Intelligent Vehicles}, 2023.

\bibitem{basile2022} G.~Basile \emph{et al.}, ``Roadmap and challenges for
reinforcement learning control in railway virtual coupling,'' \emph{Discover
Artificial Intelligence}, vol.~2, no.~27, 2022.

\bibitem{liu2025} H.~Liu \emph{et al.}, ``Coordinated operation strategy design
for virtually coupled train set: a multiagent reinforcement learning
approach,'' \emph{IEEE ITS Magazine}, vol.~17, no.~4, 2025.

\bibitem{yang2025} L.~Yang \emph{et al.}, ``Single-train multi-body MPC with
cyclic air braking and coupler-impulse minimization on the Shuozhou profile,''
in \emph{Proc.\ IEEE ITSC}, 2025.

\end{thebibliography}

\end{document}
